% ============================================================
%  10_infrastructure.tex — Section 10 : Infrastructure d'experimentation
% ============================================================
\section{Infrastructure d'expérimentation systématique}

L'exploration des contraintes CSP nécessite de tester systématiquement les combinaisons de
flags sur un même dataset (par exemple tous les benzénoïdes \texttt{h3}, \texttt{h4},
\texttt{h5} de BenzAI~DB). Trois scripts de batch ont été ajoutés dans
\texttt{csp\_solver/experiments/}.

\subsection{\texttt{batch\_main.py} — un dossier, une config}

Lance \texttt{main.py --validate} sur tous les fichiers \texttt{.graph} d'un dossier avec un
ensemble de drapeaux donnés. Génère le \texttt{data.json} de la config et régénère
\texttt{view.html} agrégé en fin d'exécution.

\begin{lstlisting}[style=benzai]
python batch_main.py plane/benzdb/h4 --validate --n-runs 10
python batch_main.py plane/benzdb/h4 --validate --no-freeze --n-runs 10
\end{lstlisting}

\subsection{\texttt{batch\_all.py} — les 8 configurations CSP}

Lance \texttt{batch\_main.py} pour chacune des $2^3 = 8$ combinaisons des trois flags
\texttt{--no-freeze}, \texttt{--no-table}, \texttt{--adj-57}.

\begin{center}
\begin{tabular}{ll}
  \toprule
  Configuration & Flags \\
  \midrule
  \texttt{default}                       & (aucun) \\
  \texttt{no-freeze}                     & \texttt{--no-freeze} \\
  \texttt{no-table}                      & \texttt{--no-table} \\
  \texttt{adj-57}                        & \texttt{--adj-57} \\
  \texttt{no-freeze\_no-table}           & \texttt{--no-freeze --no-table} \\
  \texttt{adj-57\_no-freeze}             & \texttt{--adj-57 --no-freeze} \\
  \texttt{adj-57\_no-table}              & \texttt{--adj-57 --no-table} \\
  \texttt{adj-57\_no-freeze\_no-table}   & \texttt{--adj-57 --no-freeze --no-table} \\
  \bottomrule
\end{tabular}
\end{center}

À la fin de l'exécution, \texttt{batch\_all.py} écrit un fichier \texttt{batch\_meta.json}
au niveau \texttt{output/hX/} qui contient les statistiques globales du dernier run (durée
totale, nombre d'instances, nombre de solutions trouvées, options propagées). Le viewer
agrégé affiche ce méta-fichier sous forme de bandeau si présent.

\begin{lstlisting}[style=benzai]
python batch_all.py plane/benzdb/h4 --n-runs 10
\end{lstlisting}

\subsection{Structure de sortie complète}

\begin{lstlisting}[style=benzai]
csp_solver/experiments/output/
  hX/
    batch_meta.json                 # Stats du dernier batch_all
    view.html                       # Viewer agrege (genere par view.py)
    config_name/
      data.json                     # Resultats agreges de cette config
      mol_name/
        mol_name_original.xyz       # Original (avant xTB)
        mol_name_original_opt.xyz   # Original (apres xTB)
        solutions/
          sol_K_X_Y_Z_.../          # 1 dossier par solution CSP
            source.xyz
            run_01_opt.xyz
            ...
            run_N_opt.xyz
\end{lstlisting}

\subsection{Format \texttt{data.json}}

Chaque \texttt{data.json} suit un schéma stable consommé par tous les scripts de
visualisation. Une solution typique :

\begin{lstlisting}[style=benzai]
{
  "source": "h4", "config": "default",
  "generated": "2026-04-23T15:30:00",
  "threshold_deg": 10.0, "n_runs": 10,
  "molecules": {
    "0-5-6-11": {
      "name": "0-5-6-11",
      "original": { "planar": true, "angle_deg": 0.07, "rmsd": ..., "height": ... },
      "solutions": [
        {
          "sizes": "v0=5 v1=7 v2=7 v3=5",
          "file": "sol_1_5_7_7_5/run_10_opt.xyz",
          "planar": true, "angle_deg": 0.15,    /* dernier run, retrocompat */
          "runs": {                              /* multi-runs (optionnel) */
            "n": 10, "planar_count": 9, "non_planar_count": 1,
            "planar_pct": 90, "angle_mean": 4.2, "angle_std": 1.8,
            "angle_min": 0.8, "angle_max": 18.1,
            "angles": [0.8, 1.1, ...],
            "classification": "mostly_planar"
          }
        }
      ]
    }
  }
}
\end{lstlisting}

\begin{remarque}
  Le bloc \texttt{runs} est \textbf{additif} : si absent (cas single-run), l'affichage
  retombe sur les champs classiques \texttt{planar} / \texttt{angle\_deg}. Cela garantit la
  rétrocompatibilité avec les anciens \texttt{data.json} et permet de mélanger dans le même
  rapport des tailles testées en single-run et en multi-runs.
\end{remarque}
